\renewcommand{\footerfilename}{Win1x.tex}

\section{Windows 10/11}

The capture program of choice is Demetra\footnote{https://www.shelyak.com/software/demetra/?lang=en}
on the Win10/11 platform. It allows for both capture and reduction. We focus on the image
capture part here. The main advantage is it will display the most recent image together, allow
the ds9 equivalent of a projection\index{Demetra!projection}. 

The dispersion set command: select the middle of the dispersion, then set the width.

If you double-click left-mouse, you can then start high and drag a small statistics
box over the dispersion. Cheat a bit, choose only within the bright margins of the
trace. This will return an average. A quick-and-dirty way to make sure the exposure
is good. It returns an average pixel value. Multiply that in your brain by the width
of the dispersion axis -- say 20 and you have a decent guess to the SNR for the reduced
spectrum. Adjust accordingly for roughly 10,000 x 2 = 20,000 ADU. This will assure
a strong signal.

The headers are very bad. Make sure to keep a paper logbook for each exposure. This
convention has proved workable.

\vspace{-.15cm}
\begin{enumerate}\addtolength{\itemsep}{-0.5\baselineskip}
%\setcounter{enumi}{N}
   \item   Slew to the target.
   \item   Enter the name of the target, sans spaces: HD12345
   \item   Set the number of exposures to at least 3 to combat cosmic rays, and the exptime
  to return around 20,000 adu.
   \item   After this completes, append an \_Cal to the name
   \item   Take the calibration lamp exposures (3x) 20 seconds.
   \item   Slew to the reference star
   \item   Remove the \_Cal and add the reference star: field will be HD12345\HD98765\_Ref
  This says the reference image (target name) for the previous HD science images.
   \item   Then use HD12345\HD98765\_Ref\Cal for the matching calibration images.\end{enumerate}



This helps to combat the lack of mount's RA/Dec and a Target Name.

